\chapter{Benchmarking ST against Texas Instruments}
\label{chap:vs-ti}

\section*{Introduction}

Texas Instruments (TI) is one of the world's largest semiconductor manufacturers, with a
broad portfolio of embedded processors and microcontrollers~\cite{ti-about}. Its MSPM0 and
MSPM3 families, built around Arm Cortex-M0+ and Cortex-M33 cores respectively, target
cost-sensitive and mixed-signal applications with a strong emphasis on analog
integration~\cite{ti-mspm33-product}. Unlike GigaDevice, which competes primarily on
pin-compatibility, TI competes on a fundamentally different SDK architecture: a
SysConfig-centric, unified serial peripheral model and a single-layer driver library
(driverlib) rather than the dual HAL/LL approach used by STMicroelectronics.

This chapter presents the first phase of our ST-versus-TI comparison: a comprehensive
\textbf{static analysis} of the low-level driver layers of both SDKs. Unlike the
runtime-focused benchmarks in Chapter~\ref{chap:vs-gd32}, this analysis examines source-code
quality and SDK breadth without executing firmware on hardware. We compare the
STM32CubeC5 Low-Layer (LL) drivers~\cite{stm32cubec5-sw} against the TI MSPM33 SDK
driverlib~\cite{ti-mspm33-sdk} across seven quality dimensions --- lines of code,
documentation coverage, code duplication, cyclomatic complexity, API surface design,
MISRA~C:2012 compliance, and security (CWE analysis). For the scope definition, we also
map the SDK coverage in terms of peripheral drivers, middleware, cryptographic algorithms,
and communication protocols, identifying what each vendor offers exclusively.

All measurements were produced by automated scripts using
\texttt{cloc}~\cite{cloc-tool}, \texttt{cppcheck}~\cite{cppcheck-tool} (with its
\texttt{--addon=misra} and \texttt{--addon=cert} modes),
\texttt{lizard}~\cite{lizard-tool} for cyclomatic complexity, and
\texttt{jscpd}~\cite{jscpd-tool} for duplication detection.

% ====================================================================
\section{Hardware platforms}
\label{sec:ti-boards}

The analysis in this chapter targets driver code for two boards hosting comparable
Cortex-M33 parts. Table~\ref{tab:ti-board-comparison} summarises their key characteristics.

\begin{table}[H]
\centering
\caption{Board and MCU comparison --- STM32C5A3 vs.\ MSPM33C321A.}
\label{tab:ti-board-comparison}
\begin{tabular}{l l l}
\toprule
\textbf{Feature} & \textbf{STM32C5A3 (NUCLEO-C5A3ZG)} & \textbf{MSPM33C321A (LP-MSPM33C321A)} \\
\midrule
Core & Arm Cortex-M33 (FPU, DSP, MPU) & Arm Cortex-M33 (FPU, DSP, TrustZone) \\
Max frequency & 144\,MHz & 160\,MHz \\
Flash & 1\,MB (ECC, dual-bank) & 1\,MB (ECC, dual-bank, address swap) \\
SRAM & 256\,KB (128\,KB with ECC) & 256\,KB (ECC) \\
ADC & 3 $\times$ 12-bit, 4.5\,MSPS & 2 $\times$ 12-bit, 9.4\,MSPS \\
DAC & 1 $\times$ 12-bit & 2 $\times$ 8-bit \\
CAN & 2 $\times$ FDCAN & 2 $\times$ CAN\,FD \\
USB & USB\,2.0 FS (Host/Device/DRD) & --- \\
Ethernet & Yes (MAC with DMA) & --- \\
I3C & Yes & --- \\
Debug probe & ST-LINK & XDS110 (EnergyTrace) \\
\bottomrule
\end{tabular}
\end{table}

Both parts share the same core architecture and identical flash/RAM capacities, which makes
the comparison fair at the silicon level. The STM32C5A3 offers broader connectivity (USB,
Ethernet, I3C), while the MSPM33C321A favours higher analog bandwidth (9.4\,MSPS ADC) and
provides TrustZone support~\cite{ti-mspm33-ds,stm32c5a3-product}.

% ====================================================================
\section{Static analysis of the driver layer}
\label{sec:ti-static-analysis}

This section compares the LL driver layer of STM32CubeC5 against the driverlib of the
MSPM33 SDK. The LL layer was chosen for ST because it is the closest architectural
equivalent to TI's driverlib: both are thin, inline-heavy abstractions over peripheral
registers, designed for performance-critical code. We analyse seven complementary
dimensions.

% --------------------------------------------------------------------
\subsection{Scope and methodology}
\label{subsec:ti-scope}

The analysis targets the following source sets:

\begin{itemize}
  \item \textbf{ST:} all 33 header files in
        \texttt{stm32c5xx\_drivers/ll/} (the LL layer of STM32CubeC5).
  \item \textbf{TI:} all 85 source and header files in
        \texttt{source/ti/driverlib/} and its \texttt{m33/} subdirectory
        (the driverlib of MSPM33 SDK v1.03.00.01).
\end{itemize}

Both sets were analysed on the same workstation under identical tool configurations. The
tools and their roles are:

\begin{itemize}
  \item \textbf{cloc} v2.08~\cite{cloc-tool} --- counts blank, comment, and code lines per
        language and file.
  \item \textbf{lizard}~\cite{lizard-tool} --- computes cyclomatic complexity, number of
        parameters, and token count per function.
  \item \textbf{cppcheck} v2.x~\cite{cppcheck-tool} --- runs MISRA~C:2012 addon
        (\texttt{--addon=misra}) and CWE/CERT checks.
  \item \textbf{jscpd}~\cite{jscpd-tool} --- detects copy-pasted code blocks (duplication).
  \item A custom Python script (\texttt{analyze\_firmware.py}) --- extracts API surface
        metrics (public/static/inline functions, macros, enums, naming patterns, parameter
        distributions) and documentation coverage (Doxygen blocks, \texttt{@brief},
        \texttt{@param}, \texttt{@return} tags, file headers).
\end{itemize}

% --------------------------------------------------------------------
\subsection{Lines of code}
\label{subsec:ti-loc}

Table~\ref{tab:ti-loc} presents the line-count breakdown for each SDK's driver layer.

\begin{table}[H]
\centering
\caption{Lines-of-code comparison --- ST LL vs.\ TI driverlib.}
\label{tab:ti-loc}
\begin{tabular}{l r r}
\toprule
\textbf{Metric} & \textbf{ST (LL)} & \textbf{TI (driverlib)} \\
\midrule
Source files      &        33 &        85 \\
Blank lines       &     5\,744 &     8\,747 \\
Comment lines     &    57\,649 &    45\,763 \\
Code lines        &    22\,292 &    27\,972 \\
\midrule
Total lines       &    85\,685 &    82\,482 \\
Comment-to-code ratio & 2.59 & 1.64 \\
\bottomrule
\end{tabular}
\end{table}

ST delivers its LL layer in 33 header-only files containing 22\,292 lines of code, while
TI spreads its driverlib across 85 files (both \texttt{.c} and \texttt{.h}) with
27\,972 code lines. Despite having 25\% more code, TI's total line count is slightly
lower because ST invests significantly more in inline comments: 57\,649 comment lines
versus 45\,763, yielding a comment-to-code ratio of 2.59 for ST compared with 1.64 for TI.
This difference reflects ST's header-only, fully-documented inline function strategy, where
each function carries extensive Doxygen annotations embedded alongside the implementation.

% --------------------------------------------------------------------
\subsection{Documentation coverage}
\label{subsec:ti-doc}

We measured documentation quality through Doxygen tag extraction.
Table~\ref{tab:ti-doc} summarises the results.

\begin{table}[H]
\centering
\caption{Documentation coverage --- ST LL vs.\ TI driverlib.}
\label{tab:ti-doc}
\begin{tabular}{l r r}
\toprule
\textbf{Metric} & \textbf{ST (LL)} & \textbf{TI (driverlib)} \\
\midrule
Total functions         &   3\,713 &   3\,056 \\
Documented functions    &   3\,613 &   2\,525 \\
\textbf{Doc.\ coverage} & \textbf{97.3\,\%} & \textbf{82.6\,\%} \\
\midrule
Doxygen blocks          &   5\,967 &   3\,357 \\
\texttt{@brief} tags    &   3\,826 &   5\,568 \\
\texttt{@param} tags    &   4\,440 &   3\,970 \\
\texttt{@return} tags   &   1\,705 &   2\,172 \\
File headers            &      33  &      85  \\
Inline comments         &       0  &     283  \\
\bottomrule
\end{tabular}
\end{table}

ST achieves 97.3\% function-level documentation coverage, leaving only 100 undocumented
functions. TI reaches 82.6\%, with 531 undocumented functions --- mostly in \texttt{.c}
implementation files, where TI concentrates complex logic (e.g.\ flash controller
operations, PLL configuration) without per-function Doxygen blocks. This architectural
choice is deliberate: TI documents the \emph{header} API thoroughly but omits Doxygen in
\texttt{.c} bodies, whereas ST's header-only approach means every function definition is
also its documented declaration.

Interestingly, TI uses more \texttt{@brief} tags (5\,568 vs.\ 3\,826) and more
\texttt{@return} tags (2\,172 vs.\ 1\,705), suggesting that when TI does document a
function, it tends to be descriptive about purpose and return values. ST, however, produces
78\% more Doxygen blocks overall (5\,967 vs.\ 3\,357), reflecting its higher raw volume of
documented entities.

% --------------------------------------------------------------------
\subsection{Code duplication}
\label{subsec:ti-dup}

Table~\ref{tab:ti-dup} presents the duplication metrics.

\begin{table}[H]
\centering
\caption{Code duplication --- ST LL vs.\ TI driverlib.}
\label{tab:ti-dup}
\begin{tabular}{l r r}
\toprule
\textbf{Metric} & \textbf{ST (LL)} & \textbf{TI (driverlib)} \\
\midrule
Total lines analysed     &    85\,685 &    82\,403 \\
Code lines               &    10\,461 &    18\,008 \\
Duplicate blocks         &        71 &       414 \\
Duplicate lines          &       229 &     1\,306 \\
\textbf{Duplication \%}  & \textbf{2.19\,\%} & \textbf{7.25\,\%} \\
\midrule
Same-file clones         &        71 &       338 \\
Cross-file clones        &         0 &        76 \\
\bottomrule
\end{tabular}
\end{table}

TI's driver layer exhibits 3.3$\times$ more duplication than ST's (7.25\% vs.\ 2.19\%).
All 71 duplicate blocks in ST occur within the same file, typically in symmetric
getter/setter pairs for multi-channel peripherals (e.g.\ ADC rank configuration).
TI, by contrast, has 76 cross-file clones --- code that is structurally identical across
different peripheral modules. This pattern arises from TI's UniComm architecture, where
\texttt{dl\_unicommuart.h}, \texttt{dl\_unicommspi.h}, and \texttt{dl\_unicommi2cc.h}
share a common General Serial Controller (GSC) base but replicate similar interrupt-handling
and configuration boilerplate in each specialisation.

The higher duplication in TI is not inherently a quality defect --- it reflects a design
choice that trades some code reuse for peripheral-specific readability. However, from a
maintenance perspective, ST's lower duplication means fewer locations to update when fixing
a bug or adapting to silicon errata.

% --------------------------------------------------------------------
\subsection{Cyclomatic complexity}
\label{subsec:ti-complexity}

Cyclomatic complexity (CC) measures the number of linearly independent paths through a
function; higher values indicate more decision branches and, consequently, harder-to-test
code~\cite{lizard-tool}. We ran Lizard on every function in both driver sets and
summarise the results in Table~\ref{tab:ti-complexity}.

\begin{table}[H]
\centering
\caption{Cyclomatic complexity comparison --- ST LL vs.\ TI driverlib.}
\label{tab:ti-complexity}
\begin{tabular}{l r r}
\toprule
\textbf{Metric} & \textbf{ST (LL)} & \textbf{TI (driverlib)} \\
\midrule
Analysed functions       &    3\,542 &    2\,539 \\
Mean CC                  &      4.44 &      7.35 \\
Median CC                &         4 &         5 \\
Maximum CC               &        49 &       141 \\
\midrule
Functions with CC $>$ 10 &    44 \scriptsize(1.2\,\%) &   335 \scriptsize(13.2\,\%) \\
Functions with CC $>$ 20 &    12 \scriptsize(0.3\,\%) &    97 \scriptsize(3.8\,\%) \\
\bottomrule
\end{tabular}
\end{table}

ST's driver code is markedly flatter: the average function has fewer than five decision
points, and only 1.2\,\% of functions exceed the widely accepted threshold of
CC\,=\,10~\cite{misra-c2012}. The maximum (CC\,=\,49) occurs in a temperature-calculation
helper (\texttt{LL\_ADC\_CALC\_TEMPERATURE}) that contains a long chain of calibration
look-ups.

TI's SDK shows a heavier tail: 13.2\,\% of its functions exceed CC\,=\,10, with a peak of
141 in a single function. Many of these high-complexity functions reside in the UniComm
serial-controller layer, where a large \texttt{switch} dispatches configuration variants
for UART, SPI, and I\textsuperscript{2}C through a shared code path. While this design
consolidates peripheral logic, it concentrates branching in ways that complicate unit testing
and formal verification.

From a maintainability standpoint, ST's approach of decomposing peripheral logic into many
small, low-complexity inline functions aligns better with safety-coding guidelines that
recommend CC\,$\leq$\,10 per function.

% --------------------------------------------------------------------
\subsection{API surface}
\label{subsec:ti-api}

Table~\ref{tab:ti-api} compares the API design characteristics.

\begin{table}[H]
\centering
\caption{API surface comparison --- ST LL vs.\ TI driverlib.}
\label{tab:ti-api}
\begin{tabular}{l r r}
\toprule
\textbf{Metric} & \textbf{ST (LL)} & \textbf{TI (driverlib)} \\
\midrule
Public functions         &    3\,543 &    2\,500 \\
Static functions         &         1 &        42 \\
Inline functions         &    3\,543 &    2\,032 \\
Functional macros        &       194 &        27 \\
Constant macros          &     3\,993 &    3\,271 \\
Enums                    &         1 &       361 \\
Typedefs                 &         0 &       470 \\
\midrule
Avg.\ parameters/func.  &      1.17 &      1.69 \\
Max parameters           &         8 &         8 \\
\texttt{extern "C"} guards &    33 &        43 \\
Include guards           &        33 &         0 \\
\midrule
\multicolumn{3}{l}{\textbf{Parameter distribution (functions by param count)}} \\
0 parameters             &       492 &       120 \\
1 parameter              &     2\,303 &     1\,160 \\
2 parameters             &       658 &       933 \\
3 parameters             &       135 &       274 \\
4+ parameters            &        48 &       139 \\
\bottomrule
\end{tabular}
\end{table}

The two APIs reflect fundamentally different design philosophies:

\begin{itemize}
  \item \textbf{ST LL} exposes 3\,543 public functions, \emph{all} of which are
        \texttt{\_\_STATIC\_INLINE}. The API is almost entirely composed of one-parameter
        getters and zero-parameter status queries (65\% of functions take 0 or 1 parameter).
        Type safety relies on constant macros rather than enums. The naming convention
        follows the \texttt{LL\_PERIPHERAL\_Action} pattern consistently across all 33
        files.
  \item \textbf{TI driverlib} exposes 2\,500 public functions, of which 2\,032 are inline
        (81\%). It makes heavier use of enums (361 enum types) and typedefs (470) for
        configuration parameters, which improves type safety and IDE autocompletion at the
        cost of a more complex type graph. The naming convention follows
        \texttt{DL\_Peripheral\_Action} uniformly. Functions tend to take more parameters
        on average (1.69 vs.\ 1.17), reflecting a more configuration-object-oriented
        style where a single call sets multiple related fields.
\end{itemize}

ST's 194 functional macros versus TI's 27 show that ST still uses preprocessor computation
in places (e.g.\ register-offset calculations, bit-field assembly) where TI has migrated
the same logic into inline functions. While both approaches compile to equivalent machine
code, the inline-function approach provides better debugging and type checking.

% --------------------------------------------------------------------
\subsection{MISRA~C:2012 compliance}
\label{subsec:ti-misra}

Both SDKs were analysed using \texttt{cppcheck} with the MISRA~C:2012
addon~\cite{misra-c2012}. Table~\ref{tab:ti-misra} presents the per-rule breakdown.

\begin{table}[H]
\centering
\caption{MISRA~C:2012 violations --- ST LL vs.\ TI driverlib.}
\label{tab:ti-misra}
\begin{tabular}{l l r r}
\toprule
\textbf{Rule} & \textbf{Description} & \textbf{ST} & \textbf{TI} \\
\midrule
17.3 & Function declared implicitly       & 3\,541 &  281 \\
11.4 & Conversion between pointer and integer & 84 &    4 \\
10.6 & Composite expression assigned to wider type & 29 & --- \\
20.7 & Macro parameter not enclosed in parentheses & 21 & --- \\
7.3  & Lower-case ``l'' suffix on literal  &    18 & --- \\
10.4 & Incompatible essential types in operation & --- & 67 \\
12.1 & Precedence of operators unclear     & --- &   52 \\
18.4 & Pointer arithmetic                  & --- &   76 \\
15.5 & Multiple return statements          &     4 &   14 \\
10.1 & Operands of inappropriate essential type & --- & 13 \\
3.1  & Comment contains nested comment     & --- &   11 \\
Other & Various minor rules                &    6 &   33 \\
\midrule
\textbf{Total} & & \textbf{3\,703} & \textbf{551} \\
\bottomrule
\end{tabular}
\end{table}

\subsubsection*{Interpretation}
At first glance, ST's LL layer appears to have far more MISRA violations (3\,703 vs.\ 551).
However, this requires careful interpretation:

\begin{enumerate}
  \item \textbf{Rule 17.3 dominates ST's count.} Of ST's 3\,703 violations, 3\,541 (96\%)
        are Rule~17.3 (implicit function declaration). These arise because every LL
        function is defined as \texttt{\_\_STATIC\_INLINE} in a header and \texttt{cppcheck}
        analyses each header in isolation, reporting each inline function call within the
        same header as an implicit declaration when the callee is defined later in the file.
        In a real compilation unit (where the header is \texttt{\#include}d), these are all
        resolved. This is a \emph{tooling artifact}, not a genuine code-quality issue.

  \item \textbf{Rule 11.4 is unavoidable in MCU drivers.} The 84 Rule~11.4 violations in ST
        are casts between integers and peripheral base-address pointers
        (\texttt{(ADC\_TypeDef~*)ADC1\_BASE}). Both vendors do this; TI has the same
        pattern but fewer instances (4) because TI defines peripheral pointers via
        device-header macros that \texttt{cppcheck} does not fully expand.

  \item \textbf{TI's violations are more diverse.} TI triggers 15 distinct MISRA rules
        compared with ST's 9. Rules 18.4 (pointer arithmetic, 76 occurrences), 10.4
        (essential-type mismatch, 67), and 12.1 (operator precedence, 52) point to more
        complex expression patterns in TI's \texttt{.c} implementation files.
\end{enumerate}

\noindent
If we exclude the tooling-artifact Rule~17.3 from both vendors, the adjusted counts are
ST~=~162 and TI~=~270, putting ST ahead on genuine MISRA conformance per line of code.

\subsubsection*{Security analysis (CWE)}
Both driver layers returned \textbf{zero CWE findings} from \texttt{cppcheck}'s CERT/CWE
checker. Neither SDK contains patterns flagged as common weakness enumerations, which is
expected for thin, register-level driver code that performs no dynamic allocation, string
handling, or network I/O.

% ====================================================================
\section{SDK coverage comparison}
\label{sec:ti-sdk-coverage}

Beyond driver-layer quality, the breadth of an SDK --- how many peripherals, middleware
stacks, and communication protocols it supports out of the box --- determines the
development effort required for a real product. We mapped the complete contents of both
SDKs file by file and identified exclusive features on each side.

% --------------------------------------------------------------------
\subsection{Peripheral driver coverage}
\label{subsec:ti-periph-gaps}

Table~\ref{tab:ti-periph-exclusive} lists the peripheral drivers present in only one
SDK.\footnote{For brevity, only a representative subset of ST's 25 exclusive drivers is
shown; the full list is available in the analysis output.}

\begin{table}[H]
\centering
\caption{Exclusive peripheral drivers --- features in one SDK only.}
\label{tab:ti-periph-exclusive}
\begin{tabular}{l p{5.5cm} l}
\toprule
\textbf{Vendor} & \textbf{Driver} & \textbf{Capability} \\
\midrule
\multicolumn{3}{l}{\textit{ST-exclusive (25 drivers)}} \\
ST & CORDIC accelerator              & HW trigonometric/hyperbolic math \\
ST & DAC (12-bit)                    & Full DAC peripheral \\
ST & Ethernet MAC                    & On-chip networking \\
ST & USB Host / Device / DRD         & Full USB stack support \\
ST & I3C controller                  & Next-gen I3C bus \\
ST & HASH accelerator (SHA)          & HW SHA-1/SHA-256 \\
ST & OPAMP                           & On-chip operational amplifier \\
ST & ICACHE controller               & Configurable instruction cache \\
ST & XSPI (Octo/Quad-SPI)           & Advanced memory-mapped SPI \\
ST & LPUART, USART, SmartCard, SMBus & Extended serial protocols \\
\midrule
\multicolumn{3}{l}{\textit{TI-exclusive (16 drivers)}} \\
TI & High-Speed ADC (HSADC)          & Dedicated 9.4\,MSPS ADC block \\
TI & General Serial Controller (GSC) & Unified serial engine \\
TI & UniComm (I2C/SPI/UART)          & Configurable serial block \\
TI & Signal Processing Subsystem     & On-chip analog processing \\
TI & Low-Frequency Subsystem (LFSS)  & Unified RTC + tamper + scratchpad \\
TI & Key Store Controller            & HW key management \\
TI & Error Action Module             & HW automatic fault handling \\
TI & UART extend (IrDA, Manchester)  & Extended UART modes \\
\bottomrule
\end{tabular}
\end{table}

ST leads with 25 exclusive peripheral drivers to TI's 16. The most significant gaps are
in connectivity: ST offers USB, Ethernet, and I3C, which are entirely absent from the
MSPM33C321A. TI compensates with stronger analog integration (high-speed ADC, signal
processing subsystem) and a unified serial architecture that consolidates UART, SPI, and
I2C into a single configurable peripheral.

% --------------------------------------------------------------------
\subsection{Middleware coverage}
\label{subsec:ti-mw-gaps}

Table~\ref{tab:ti-mw-exclusive} compares the middleware stacks.

\begin{table}[H]
\centering
\caption{Exclusive middleware --- present in one SDK only.}
\label{tab:ti-mw-exclusive}
\begin{tabular}{l l p{6cm}}
\toprule
\textbf{Vendor} & \textbf{Middleware} & \textbf{Purpose} \\
\midrule
ST & FileX              & FAT-compatible embedded file system \\
ST & LevelX             & Flash wear-leveling (NAND/NOR) \\
ST & LwIP               & Lightweight TCP/IP stack \\
ST & USBX               & USB device/host class drivers \\
ST & mbedTLS             & TLS/SSL library (TLS~1.2/1.3) \\
ST & Open Bootloader     & Multi-interface bootloader \\
ST & stFCF               & Security provisioning framework \\
ST & stCryptoLib          & Comprehensive crypto library \\
\midrule
TI & Edge AI             & On-device ML inference \\
TI & LVGL                & Embedded graphics/GUI library \\
TI & LIN Bus             & Automotive LIN protocol \\
TI & DALI                & Smart lighting protocol \\
TI & IQmath              & Fixed-point DSP math library \\
TI & Post-Quantum Crypto & ML-DSA (Dilithium) signatures \\
TI & Curve25519          & Modern elliptic curve crypto \\
\bottomrule
\end{tabular}
\end{table}

ST provides 8 exclusive middleware components focused on connectivity (USB, Ethernet/IP,
file system) and security (TLS, crypto library, secure boot). TI provides 7 exclusive
components that target emerging domains: Edge~AI for on-device machine learning, LVGL for
embedded GUIs, automotive protocols (LIN, DALI), and forward-looking cryptography
(post-quantum ML-DSA, Curve25519).

% --------------------------------------------------------------------
\subsection{Cryptographic algorithm coverage}
\label{subsec:ti-crypto}

Table~\ref{tab:ti-crypto} provides a detailed comparison of the algorithms available in
each vendor's crypto libraries.

\begin{table}[H]
\centering
\caption{Cryptographic algorithm coverage comparison.}
\label{tab:ti-crypto}
\begin{tabular}{l c c}
\toprule
\textbf{Algorithm} & \textbf{ST (stCryptoLib)} & \textbf{TI (crypto/)} \\
\midrule
AES (ECB, CBC, CTR, GCM)  & \checkmark & \checkmark \\
AES-CCM / AES-XTS         & \checkmark & ---        \\
SHA-224/256                & \checkmark & \checkmark \\
SHA-384/512, SHA-3         & \checkmark & ---        \\
HMAC                       & \checkmark & \checkmark \\
CMAC / KMAC                & \checkmark & ---        \\
RSA (PKCS\#1)              & \checkmark & \checkmark \\
ECDSA / ECDH               & \checkmark & \checkmark \\
EdDSA (Ed25519)            & \checkmark & \checkmark \\
SM2 / SM3 (Chinese std.)   & \checkmark & ---        \\
ChaCha20-Poly1305          & \checkmark & ---        \\
ML-DSA (Post-Quantum)      & ---        & \checkmark \\
Curve25519 / X25519         & ---        & \checkmark \\
TRNG (HW)                  & \checkmark & \checkmark \\
PKA (HW)                   & \checkmark & \checkmark \\
Key Store (HW)              & \checkmark & \checkmark \\
\bottomrule
\end{tabular}
\end{table}

ST offers the broader algorithm portfolio (10+ exclusive algorithms), covering legacy
standards (SHA-1), Chinese national standards (SM2/SM3), and authenticated-encryption
modes (AES-CCM, ChaCha20-Poly1305). TI, while narrower, holds a forward-looking advantage
with post-quantum cryptography (ML-DSA/Dilithium) and modern elliptic-curve primitives
(Curve25519), which positions it well for long-term security requirements~\cite{nist-pqc}.

% --------------------------------------------------------------------
\subsection{Communication protocol coverage}
\label{subsec:ti-comm}

Table~\ref{tab:ti-comm-protocols} summarises the communication protocol support.

\begin{table}[H]
\centering
\caption{Communication protocol support comparison.}
\label{tab:ti-comm-protocols}
\begin{tabular}{l c c}
\toprule
\textbf{Protocol} & \textbf{ST} & \textbf{TI} \\
\midrule
UART / SPI / I2C / I2S / CAN\,FD   & \checkmark & \checkmark \\
USB (Device + Host)                  & \checkmark & ---        \\
Ethernet                             & \checkmark & ---        \\
I3C                                  & \checkmark & ---        \\
SMBus / SmartCard (ISO~7816)         & \checkmark & ---        \\
LPUART / USART (synchronous)         & \checkmark & ---        \\
IrDA mode / Manchester encoding      & ---        & \checkmark \\
LIN Bus                              & ---        & \checkmark \\
\bottomrule
\end{tabular}
\end{table}

ST supports 7 exclusive communication protocols, with USB and Ethernet being the most
impactful for connected applications. TI brings 3 exclusive capabilities: IrDA mode,
Manchester encoding (both through its extended UART), and LIN Bus for automotive
applications.

% ====================================================================
\section{Summary of findings}
\label{sec:ti-static-summary}

Table~\ref{tab:ti-summary-scoreboard} provides a consolidated view of the static-analysis
results.

\begin{table}[H]
\centering
\caption{Static-analysis scoreboard --- STM32CubeC5 LL vs.\ TI MSPM33 driverlib.}
\label{tab:ti-summary-scoreboard}
\begin{tabular}{l l l l}
\toprule
\textbf{Dimension} & \textbf{ST result} & \textbf{TI result} & \textbf{Advantage} \\
\midrule
Code lines            & 22\,292          & 27\,972          & ST (leaner) \\
Documentation coverage & 97.3\,\%        & 82.6\,\%         & ST \\
Code duplication       & 2.19\,\%        & 7.25\,\%         & ST \\
Cyclomatic complexity  & 4.44 avg        & 7.35 avg         & ST (flatter) \\
API functions          & 3\,543          & 2\,500           & ST (wider) \\
Type safety (enums)    & 1               & 361              & TI \\
MISRA violations       & 162\textsuperscript{*} & 270       & ST \\
Security (CWE)         & 0               & 0                & Tie \\
Exclusive peripherals  & 25              & 16               & ST \\
Exclusive middleware   & 8               & 7                & Balanced \\
Crypto algorithms      & 10+ exclusive   & 2 exclusive      & ST (breadth), TI (PQC) \\
Comm.\ protocols       & 7 exclusive     & 3 exclusive      & ST \\
\bottomrule
\end{tabular}
\par\smallskip
{\footnotesize \textsuperscript{*}After excluding tooling-artifact Rule~17.3 from both vendors.}
\end{table}

The static analysis reveals that ST's LL driver layer is leaner, better documented, less
duplicated, and more MISRA-conformant when tooling artifacts are factored out. ST also holds
a clear lead in SDK breadth, particularly in connectivity (USB, Ethernet, I3C) and security
middleware (TLS, comprehensive crypto library).

TI, however, brings distinct strengths: stronger type safety through enums and typedefs,
forward-looking post-quantum cryptography, superior analog integration (9.4\,MSPS ADC), a
unified serial architecture that simplifies board design, and emerging-domain middleware
(Edge~AI, LVGL, automotive protocols). These advantages position TI well for mixed-signal
and AI-at-the-edge applications where connectivity is less critical than analog performance.

\newpage
\section*{Conclusion}

This chapter presented the first phase of our ST-versus-TI comparison, focusing on a
comprehensive static analysis of the low-level driver layers of STM32CubeC5 and the TI
MSPM33 SDK. The analysis covered seven quality dimensions and a complete SDK coverage
mapping.

The results demonstrate that ST's ecosystem is the more mature and broadly equipped
platform, excelling in documentation discipline, code hygiene, MISRA conformance, and
middleware breadth. TI offers a credible alternative with targeted advantages in analog
performance, type-safe API design, forward-looking cryptography, and emerging middleware for
AI and automotive applications.

Future work will extend this comparison to runtime benchmarks --- footprint analysis,
execution time, and power consumption on identical workloads --- to complement the
static-analysis findings presented here with empirical measurements from both platforms.